\def\nn {neural network}
\def\nns {neural networks}
\def\Nn {Neural network}
\def\NN {Neural Network}
\def\neu {neuron}
\def\neus {neurons}

\part {Synopsis} {}

\chap Synopsis

This thesis is concerned with the implementation of, and
experimentation with, back-propagation \nn s.

Basic \nn\ models are explained, and the back-propagation model is
described in detail in an algorithmic form.

The specification, design, and performance of the constructed software 
package, fishNET, are described.
Equations for memory usage and approximate calculation and learning
times are given.

Experimentation with learning and error rates shows that learning
parameter and network parameter variations may have a large effect
on learning time, but only a small effect on error rate.
It is shown that networks which take longer to learn have a
slightly lower error rate.

Introduction of casualties in both connections and \neus\ shows
that knowledge in the network is distributed amongst the weights but
localised in \neus, which act as feature extractors.

Finally, a complete printout of fishNET is included in the
appendices, along with the dot-matrix encoded character data.

